\section*{Context and Motivation}

In the digital age, the ability to automatically extract text from images has become
increasingly important across numerous domains. \gls{ocr} technology enables the conversion
of printed or handwritten text from images into machine-readable format, facilitating
document digitization, automated data entry, text search in scanned documents, and
accessibility features for visually impaired users.

Traditional \gls{ocr} systems relied heavily on hand-crafted features and rule-based
approaches, requiring extensive domain expertise and manual tuning for different fonts,
languages, and image conditions. The advent of deep learning has revolutionized this
field, enabling end-to-end learning systems that automatically discover optimal feature
representations directly from data.

\section*{Project Objectives}

This project aims to develop a complete \gls{ocr} system for printed character recognition
using modern deep learning techniques. The specific objectives are:

\begin{enumerate}
\item Design and implement a robust preprocessing pipeline to handle varied image
conditions including different lighting, contrast, and noise levels.

\item Build a deep neural network architecture combining convolutional and recurrent
layers to extract visual features and model sequential dependencies in text.

\item Train the model on a comprehensive dataset of 422,400 character images spanning
66 classes (digits 0-9, uppercase A-Z, lowercase a-z, and common punctuation marks).

\item Implement \gls{ctc} loss for sequence-to-sequence learning without requiring
explicit character-level alignment during training.

\item Deploy the trained model as a web application with a user-friendly interface
for real-time character and text recognition.

\item Evaluate model performance and identify areas for future improvement.
\end{enumerate}

\section*{Report Structure}

This report is organized as follows:

\textbf{Chapter 1} presents the problem statement, reviews related work in \gls{ocr}
and deep learning for text recognition, and discusses the theoretical foundations of
\gls{cnn}, \gls{rnn}, \gls{lstm}, and \gls{ctc} loss.

\textbf{Chapter 2} describes the dataset structure, preprocessing techniques including
binarization and normalization, and the data loading pipeline implemented for efficient
training.

\textbf{Chapter 3} details the model architecture, including the \gls{cnn} feature
extractor, \gls{bilstm} sequence encoder, and \gls{ctc} decoder. Implementation details
using TensorFlow and Keras are provided.

\textbf{Chapter 4} presents the training procedure, hyperparameter configuration,
deployment as a Flask web application, and evaluation results with discussion of
observed behaviors and limitations.

Finally, the \textbf{Conclusion} summarizes the achievements and proposes perspectives
for future enhancements.
